% GENERATED FILE. Edit canonical lessons, then run npm run generate:books.
% canonical-source-hash: fnv1a64:409976a45433f609
% canonical-lessons: ES-C669-plano, ES-C669-orientarse, ES-C669-ubicacion, ES-C669-senalar, ES-C669-situado

\chapter{Con el plano}
\label{ch:es-con-el-plano}
\hlchaptermodality{669}

\begin{chapteropening}
\textbf{By the end of this chapter:} \emph{I can find a place on a street map, point it out and say where it is.}

Complete the last lesson of chapter 669: I can find a place on a street map, point it out and say where it is.
\end{chapteropening}

\section[el plano]{el plano — the street map, the plan}

\label{lesson:ES-C669-plano}



\begin{quote}
Before the new one: say the Spanish for the excess, then the Spanish for to drag.
\end{quote}

\subsection*{You'll want to know: el plano}
\textbf{el plano} — \textquotedblleft{}the street map, the plan\textquotedblright{}.

\begin{quote}
¿Tiene un plano de la ciudad?
\end{quote}

\begin{quote}
\emph{Do you have a map of the town?}
\end{quote}

\emph{El plano} is of a town or a building; \emph{el mapa} is of a country.

\begin{cousinweb}
\emph{plano} is a learned borrowing from Latin \emph{plānus} \textquotedblleft{}even, level, flat\textquotedblright{} (Lewis \& Short; EtymDB): a town drawn flat on paper. The inherited form is \emph{llano} \textquotedblleft{}flat\textquotedblright{}.
\end{cousinweb}

\begin{grammarlens}[title={What you've built}]
One more word for telling, explaining and describing at B1.
\end{grammarlens}

\subsection*{Your turn}
\begin{itemize}
\raggedright
  \item \emph{Say it:} \emph{el plano}
  \item \emph{Say it:} \emph{el plano}, once more
  \item \emph{Say it:} a sentence of your own with \emph{el plano}
  \item \emph{Recall:} say the Spanish for the excess, then the Spanish for to drag, then say \emph{el plano} again
\end{itemize}

\begin{tcolorbox}[breakable,colback=teal!4,colframe=teal!35!black,title={Before you move on}]
What does el plano mean? (\textquotedblleft{}The street map, the plan\textquotedblright{}.) Say it once more.
\end{tcolorbox}
\section[orientarse]{orientarse — to find your way, to get your bearings}

\label{lesson:ES-C669-orientarse}



\begin{quote}
Before the new one: say the Spanish for to drag, then the Spanish for the street map.
\end{quote}

\subsection*{You'll want to know: orientarse}
\textbf{orientarse} — \textquotedblleft{}to find your way, to get your bearings\textquotedblright{}.

\begin{quote}
Me oriento bien en las ciudades.
\end{quote}

\begin{quote}
\emph{I find my way around cities easily.}
\end{quote}

The opposite is \emph{perderse}.

\begin{cousinweb}
\emph{orientarse} is built on \emph{oriente} \textquotedblleft{}east\textquotedblright{}, a learned borrowing from Latin \emph{oriēns} \textquotedblleft{}rising\textquotedblright{}, said of the sun (EtymDB): to find your way was first to find the east.
\end{cousinweb}

\begin{grammarlens}[title={What you've built}]
One more word for telling, explaining and describing at B1.
\end{grammarlens}

\subsection*{Your turn}
\begin{itemize}
\raggedright
  \item \emph{Say it:} \emph{orientarse}
  \item \emph{Say it:} \emph{orientarse}, once more
  \item \emph{Say it:} a sentence of your own with \emph{orientarse}
  \item \emph{Recall:} say the Spanish for to drag, then the Spanish for the street map, then say \emph{orientarse} again
\end{itemize}

\begin{tcolorbox}[breakable,colback=teal!4,colframe=teal!35!black,title={Before you move on}]
What does orientarse mean? (\textquotedblleft{}To find your way, to get your bearings\textquotedblright{}.) Say it once more.
\end{tcolorbox}
\section[la ubicación]{la ubicación — the location}

\label{lesson:ES-C669-ubicacion}



\begin{quote}
Before the new one: say the Spanish for the street map, then the Spanish for to find your way.
\end{quote}

\subsection*{You'll want to know: la ubicación}
\textbf{la ubicación} — \textquotedblleft{}the location\textquotedblright{}.

\begin{quote}
Te mando la ubicación del hotel.
\end{quote}

\begin{quote}
\emph{I will send you the location of the hotel.}
\end{quote}

Latin America often uses \emph{ubicar} for to find or to be located.

\begin{cousinweb}
\emph{ubicación} is built on \emph{ubicar} \textquotedblleft{}to place, to locate\textquotedblright{}, which is built on Latin \emph{ubī} \textquotedblleft{}where\textquotedblright{} (EtymDB).
\end{cousinweb}

\begin{grammarlens}[title={What you've built}]
One more word for telling, explaining and describing at B1.
\end{grammarlens}

\subsection*{Your turn}
\begin{itemize}
\raggedright
  \item \emph{Say it:} \emph{la ubicación}
  \item \emph{Say it:} \emph{la ubicación}, once more
  \item \emph{Say it:} a sentence of your own with \emph{la ubicación}
  \item \emph{Recall:} say the Spanish for the street map, then the Spanish for to find your way, then say \emph{la ubicación} again
\end{itemize}

\begin{tcolorbox}[breakable,colback=teal!4,colframe=teal!35!black,title={Before you move on}]
What does la ubicación mean? (\textquotedblleft{}The location\textquotedblright{}.) Say it once more.
\end{tcolorbox}
\section[señalar]{señalar — to point to, to point out; to mark}

\label{lesson:ES-C669-senalar}



\begin{quote}
Before the new one: say the Spanish for to find your way, then the Spanish for the location.
\end{quote}

\subsection*{You'll want to know: señalar}
\textbf{señalar} — \textquotedblleft{}to point to, to point out; to mark\textquotedblright{}.

\begin{quote}
Me señaló el hotel en el plano.
\end{quote}

\begin{quote}
\emph{He pointed out the hotel on the map.}
\end{quote}

\emph{Señalar con el dedo} — to point with your finger.

\begin{cousinweb}
\emph{señalar} is built on \emph{señal} \textquotedblleft{}sign\textquotedblright{} (EtymDB), which is inherited from Late Latin \emph{signāle}, from \emph{signum} \textquotedblleft{}a mark, a sign\textquotedblright{}.
\end{cousinweb}

\begin{grammarlens}[title={What you've built}]
One more word for telling, explaining and describing at B1.
\end{grammarlens}

\subsection*{Your turn}
\begin{itemize}
\raggedright
  \item \emph{Say it:} \emph{señalar}
  \item \emph{Say it:} \emph{señalar}, once more
  \item \emph{Say it:} a sentence of your own with \emph{señalar}
  \item \emph{Recall:} say the Spanish for to find your way, then the Spanish for the location, then say \emph{señalar} again
\end{itemize}

\begin{tcolorbox}[breakable,colback=teal!4,colframe=teal!35!black,title={Before you move on}]
What does señalar mean? (\textquotedblleft{}To point to, to point out; to mark\textquotedblright{}.) Say it once more.
\end{tcolorbox}
\section[situado]{situado — located, situated}

\label{lesson:ES-C669-situado}



\begin{quote}
Before the new one: say the Spanish for the location, then the Spanish for to point to.
\end{quote}

\subsection*{You'll want to know: situado}
\textbf{situado} — \textquotedblleft{}located, situated\textquotedblright{}.

\begin{quote}
El hotel está situado en el centro.
\end{quote}

\begin{quote}
\emph{The hotel is located in the centre.}
\end{quote}

Used with \emph{estar}. Latin America often says \emph{ubicado}.

\begin{cousinweb}
\emph{situado} is the past form of \emph{situar} \textquotedblleft{}to place\textquotedblright{}, which is built on Latin \emph{situs} \textquotedblleft{}position, site\textquotedblright{}.
\end{cousinweb}

\begin{grammarlens}[title={What you've built}]
One more word for telling, explaining and describing at B1.
\end{grammarlens}

\subsection*{Your turn}
\begin{itemize}
\raggedright
  \item \emph{Say it:} \emph{situado}
  \item \emph{Say it:} \emph{situado}, once more
  \item \emph{Say it:} a sentence of your own with \emph{situado}
  \item \emph{Recall:} say the Spanish for the location, then the Spanish for to point to, then say \emph{situado} again
\end{itemize}

\begin{tcolorbox}[breakable,colback=teal!4,colframe=teal!35!black,title={Before you move on}]
What does situado mean? (\textquotedblleft{}Located, situated\textquotedblright{}.) Say it once more.
\end{tcolorbox}
